\section{Đá vôi, silicate và tài nguyên khoáng sản}

\subsection{Đá vôi và các sản phẩm từ đá vôi}

Đá vôi có thành phần chính là \ce{CaCO3}. Ba phản ứng quen thuộc:
\[
\ce{CaCO3 ->[t^\circ] CaO + CO2},
\]
\[
\ce{CaO + H2O -> Ca(OH)2},
\]
\[
\ce{Ca(OH)2 + CO2 -> CaCO3 v + H2O}.
\]
Các phản ứng này liên hệ trực tiếp với sản xuất vôi, xử lí môi trường và vật liệu xây dựng.

\begin{vidu}
Nung \SI{1.00}{\ton} đá vôi chứa $90\%$ \ce{CaCO3}, hiệu suất phân hủy đạt $85\%$. Tính khối lượng \ce{CaO} thu được.
\end{vidu}
\begin{loigiai}
Khối lượng \ce{CaCO3} tinh khiết:
\[
m_{\ce{CaCO3}}=1.00\times0.90=\SI{0.900}{\ton}.
\]
\bangmol{5}
{\ce{CaCO3} & \rightarrow & \ce{CaO} & + & \ce{CO2}}
{100 & & 56 & &}
Khối lượng \ce{CaO} theo lí thuyết:
\[
m_{\ce{CaO},\text{lt}}=0.900\times\frac{56}{100}=\SI{0.504}{\ton}.
\]
Khối lượng thực tế:
\[
m_{\ce{CaO}}=0.504\times0.85\approx\SI{0.428}{\ton}.
\]

\textbf{Vậy,} thu được khoảng \SI{0.428}{\ton} \ce{CaO}.
\end{loigiai}

\subsection{Silicate và vật liệu silicate}

Silicon dioxide \ce{SiO2} có nhiều trong cát và là nguyên liệu quan trọng của công nghiệp silicate. Những vật liệu quen thuộc gồm thủy tinh, gốm và xi măng. Ở mức chương trình, cần nắm nguyên liệu chính, vai trò và ứng dụng; không đi sâu cấu trúc mạng tinh thể.

\subsection{Khai thác và sử dụng tài nguyên hợp lí}

Khai thác quặng và sản xuất kim loại tiêu tốn tài nguyên, năng lượng và có thể tạo chất thải. Khi đánh giá một quy trình, hãy xem xét đồng thời:
\begin{enumerate}
\item lượng quặng cần khai thác;
\item hiệu suất thu hồi kim loại;
\item sản phẩm phụ có thể tận dụng;
\item khí thải hoặc chất thải rắn;
\item khả năng tái chế kim loại.
\end{enumerate}

\begin{notechem}
Tái chế kim loại giúp giảm nhu cầu khai thác quặng mới và thường tiết kiệm tài nguyên. Tuy nhiên, việc tái chế vẫn cần thu gom, phân loại và xử lí phù hợp.
\end{notechem}
